% エントリ 2(2026-11-03): 自由場の超放射と 2 世代目の BH
\section{[E2, 2026-11-03] 自由場の超放射: 2 世代目の BH はどこで回転を失うか}\label{sec:e02}

自己相互作用を入れる前の基準として、自由場のスカラーで、生まれたときのスピンが $\chi_0=0.69$ の BH(2 世代目)がいつ、どこまで回転を失うかを計算した。計算は \texttt{calc/sr.py}(共通部品)と \texttt{calc/e02\_freefield.py}、出力は \texttt{calc/out/e02.json}。

\subsection{式}

$\alpha\equiv GM\mu/\hbar c$ とし、$G=c=\hbar=1$ で書く。準位 $|n\ell m\rangle$ の振動数を $\omega_{n\ell m}=E_{n\ell m}+i\Gamma_{n\ell m}$ と書く(\arxiv{1908.10370} の式 2.27)。

\paragraph{エネルギー}
\arxiv{1908.10370} の式 2.28--2.29 を使う。
\begin{equation}
\frac{E_{n\ell m}}{\mu}=1-\frac{\alpha^2}{2n^2}-\frac{\alpha^4}{8n^4}+\frac{f_{n\ell}}{n^3}\alpha^4+\frac{h_\ell}{n^3}\chi m\alpha^5,
\qquad f_{n\ell}=-\frac{6}{2\ell+1}+\frac{2}{n},\quad h_\ell=\frac{16}{2\ell(2\ell+1)(2\ell+2)}.
\end{equation}

\paragraph{成長率}
\arxiv{1908.10370} の式 2.32--2.34 を使う。
\begin{equation}
\Gamma_{n\ell m}=2\tilde r_+C_{n\ell}\,g_{\ell m}\,(m\Omega_H-E_{n\ell m})\,\alpha^{4\ell+5},
\qquad
C_{n\ell}=\frac{2^{4\ell+1}(n+\ell)!}{n^{2\ell+4}(n-\ell-1)!}\Bigl[\frac{\ell!}{(2\ell)!(2\ell+1)!}\Bigr]^2,
\label{eq:e02_rate}
\end{equation}
記号は次のとおり($M=1$ の単位)。
\begin{itemize}
\item $g_{\ell m}=\prod_{k=1}^{\ell}\bigl[k^2(1-\chi^2)+(\chi m-2r_+E)^2\bigr]$。
\item $\tilde r_+=1+\sqrt{1-\chi^2}$、$\Omega_H=\chi/(2r_+)$。
\end{itemize}

\paragraph{規約の注意}
$\Gamma$ は場の振幅の成長率である。占有数 $N$ は $\mathrm{d}N/\mathrm{d}t=2\Gamma N$ で育つ。
\begin{itemize}
\item 小さい $\alpha$ では、式 \eqref{eq:e02_rate} は $\Gamma_{211}\to\mu\alpha^8(\chi-2\alpha\tilde r_+)/48$ になる。占有数の率はその 2 倍で、分母は 24 になる。
\item 本文で照合した \arxiv{1908.10370} は、式 2.32 を $\alpha\lesssim0.2$ で数値解とよく合うとしている。これより大きい $\alpha$ では相対論的な率に替える(E3 以降)。
\end{itemize}

\paragraph{BH の変化}
雲の量子 1 個が BH から $m\hbar$ の角運動量と $\hbar E$ のエネルギーを持ち去る。準位 $|n\ell m\rangle$ は $m\Omega_H(\chi)=E_{n\ell m}$ となる臨界スピンで成長を止める。

\subsection{解き方: 飽和で区切る半解析}

複数の準位を ODE で同時に解くと、飽和したあとの平衡が非常に硬く、発散した(\texttt{DESIGN.md} D3)。そこで次の半解析を使う。
\begin{itemize}
\item 各準位は 1 個の量子から育つ。$e$ 倍の回数 $A_\ell=\int2\Gamma_\ell\,\mathrm{d}t$ が $\ln N_\mathrm{sat}$ に届いたとき飽和する。
\item 飽和すると、BH は臨界スピンまで回転を失う。$N_\mathrm{sat}$ は、臨界スピンまでに渡す角運動量を $m\hbar$ で割った数である。
\item 飽和にかかる最後の数 $e$-fold は無視する($\ln N_\mathrm{sat}\approx180$ に対して数個)。
\item 準位は $211$、$322$、$433$ を入れた。
\end{itemize}
$211$ だけの場合に ODE と比べた(表~\ref{tab:e02_check})。最終スピンの差は最大 $1.4\times10^{-4}$($\alpha=0.15$)である。スピンが中点まで落ちる時刻は、時間の格子の刻み(0.7\%)以内で一致した。

\begin{table}[H]
\centering\small
\caption{半解析と ODE の照合($211$ だけ、$\chi_0=0.69$)。$t_\mathrm{mid}$ はスピンが初期値と最終値の中点まで落ちる時刻。}
\label{tab:e02_check}
\begin{tabular}{@{}rrrcccc@{}}
\toprule
$M\ [M_\odot]$ & $\mu\ [\mathrm{eV}]$ & $\alpha$ & $\chi_f$(ODE) & $\chi_f$(半解析) & $t_\mathrm{mid}$(ODE) & $t_\mathrm{mid}$(半解析) \\
\midrule
15 & $7\times10^{-13}$ & 0.079 & 0.29705 & 0.29705 & 301 yr & 301 yr \\
5 & $1\times10^{-12}$ & 0.037 & 0.14578 & 0.14578 & $4.93\times10^4$ yr & $4.93\times10^4$ yr \\
30 & $3\times10^{-13}$ & 0.067 & 0.25674 & 0.25674 & 2128 yr & 2128 yr \\
60 & $1\times10^{-13}$ & 0.045 & 0.17399 & 0.17399 & $1.28\times10^5$ yr & $1.28\times10^5$ yr \\
10 & $2\times10^{-12}$ & 0.150 & 0.53555 & 0.53569 & 1.83 yr & 1.82 yr \\
\bottomrule
\end{tabular}
\end{table}

\subsection{結果}

\begin{figure}[H]
\centering
% E2: 自由場で 2 世代目の BH(chi0 = 0.69)が回転を失う範囲。データは calc/e02_freefield.py
\pgfplotsset{colormap={spinmap}{color(0)=(cA!85!black) color(0.35)=(cA) color(0.6)=(cA!35) color(0.69)=(white)}}
\begin{tikzpicture}
\begin{groupplot}[group style={group size=2 by 1, horizontal sep=22mm},
  width=0.45\linewidth, height=0.42\linewidth, note]
% (a) t = 1e8 yr での chi_f の地図
\nextgroupplot[xmode=log, ymode=log, xmin=3, xmax=150, ymin=1e-14, ymax=1e-11,
  xlabel={$M\ [M_\odot]$}, ylabel={$\mu\ [\mathrm{eV}]$}, grid=none, axis on top,
  title={\small (a) $t=10^8\yr$ 後の $\chi_f$},
  colorbar, colormap name=spinmap, point meta min=0, point meta max=0.69,
  colorbar style={width=2.5mm, ytick={0,0.2,0.4,0.6}, yticklabel style={font=\scriptsize}, ylabel={$\chi_f$}, ylabel style={font=\scriptsize}}]
  \addplot[matrix plot*, mesh/cols=121, mesh/ordering=x varies, point meta=\thisrow{chif}]
    table[x expr={10^\thisrow{lgM}}, y expr={10^\thisrow{lgmu}}] {figs/data/e02_chif_t8.dat};
  \draw[ink2, line width=0.5pt] (axis cs:10,1e-14) -- (axis cs:10,1e-11);
  \draw[ink2, line width=0.5pt] (axis cs:20,1e-14) -- (axis cs:20,1e-11);
  \draw[ink2, line width=0.5pt, dashed] (axis cs:45,1e-14) -- (axis cs:45,1e-11);
  \addplot[cB, line width=0.8pt, densely dashed, domain=3:150, samples=2] {7.1e-13};
  \node[font=\scriptsize, color=cB, fill=white, fill opacity=0.8, text opacity=1, inner sep=1pt, anchor=south west] at (axis cs:3.3,8e-13) {$7.1\times10^{-13}$\,eV};
% (b) mu = 7.1e-13 eV での最終スピンと GW241011
\nextgroupplot[xmode=log, xmin=3, xmax=150, ymin=0, ymax=0.9,
  xlabel={$M\ [M_\odot]$}, ylabel={$\chi_f$},
  title={\small (b) $\mu=7.1\times10^{-13}$\,eV},
  legend style={font=\scriptsize, at={(0.98,0.02)}, anchor=south east, fill=white, fill opacity=0.85, text opacity=1, draw=none}, legend cell align=left]
  \fill[cM!18] (axis cs:10,0) rectangle (axis cs:20,0.9);
  \fill[cM!18] (axis cs:45,0) rectangle (axis cs:100,0.9);
  \addplot[cA!45, line width=0.9pt] table[x=M, y=chif_t6] {figs/data/e02_regge_mu7.1e-13.dat};
  \addlegendentry{$t=10^6\yr$}
  \addplot[cA, line width=0.9pt] table[x=M, y=chif_t8] {figs/data/e02_regge_mu7.1e-13.dat};
  \addlegendentry{$10^8\yr$}
  \addplot[cA!60!black, line width=0.9pt, dashed] table[x=M, y=chif_t10] {figs/data/e02_regge_mu7.1e-13.dat};
  \addlegendentry{$10^{10}\yr$}
  \addplot[ink2, line width=0.5pt, densely dotted, domain=3:150, samples=2, forget plot] {0.69};
  % GW241011 の主星(LVK 文献レビュー T04-28: m1 = 19.6 +3.6/-2.5, chi1 = 0.78 +/- 0.09、90%)
  \addplot[cB, only marks, mark=*, mark size=1.8pt, error bars/.cd, x dir=both, x explicit, y dir=both, y explicit]
    coordinates {(19.6,0.78) +- (0,0)};
  \addlegendentry{GW241011 主星}
  \draw[cB, line width=0.7pt] (axis cs:17.1,0.78) -- (axis cs:23.2,0.78);
  \draw[cB, line width=0.7pt] (axis cs:19.6,0.69) -- (axis cs:19.6,0.87);
\end{groupplot}
\end{tikzpicture}

\caption{自由場で、$\chi_0=0.69$ の BH が回転を失う範囲。灰色の帯は、GWTC-5 で速く回る部分集団が見つかった質量 $10$--$20\Msun$ と $45\Msun$ 以上(\arxiv{2605.27226})。
(a) $t=10^8\yr$ 後のスピン $\chi_f$(白は回転を失わない)。縦線は質量 10・20・$45\Msun$、橙の破線は \arxiv{2609.02678} の質量。準位 211・322・433 が飽和する範囲が斜めの帯になる($\alpha\propto M\mu$)。
(b) $\mu=7.1\times10^{-13}\eV$ での最終スピン。点は GW241011 の主星(90\%、\review\ T04-28)。}
\label{fig:e02}
\end{figure}

\begin{result}
\begin{itemize}
\item \textbf{$15\Msun$、$\mu=7\times10^{-13}\eV$}: $\alpha=0.079$、占有数の $e$-fold は $1.7\yr$、$\ln N_\mathrm{sat}=179$ である。スピンは約 300 年で中点まで落ち、$0.297$ に止まる。テーマ選定のときの見積もり(e-fold 約 1 年、育ちきるまで $\sim10^2$ 年)は桁で合っていた。
\item \textbf{GW241011 の主星}(\review\ T04-28 の $m_1=19.6\Msun$、$\chi_1=0.78$): $\mu=7.1\times10^{-13}\eV$ の雲があると、スピンは約 35 年で 90\% 区間の下端 0.69 を下回る。$10^{10}\yr$ 後は $322$ も飽和して 0.195 になる。
\item \textbf{帯の中のすべての質量で $\chi_f<0.5$ になる $\mu$}(図~\ref{fig:e02}(a)):
  \begin{center}
  \begin{tabular}{@{}lcc@{}}
  \toprule
  $t$ & $10$--$20\Msun$ & $45$--$100\Msun$ \\
  \midrule
  $10^6\yr$ & $(4.0\text{--}8.9)\times10^{-13}\eV$ & $(1.0\text{--}1.9)\times10^{-13}\eV$ \\
  $10^8\yr$ & $2.4\times10^{-13}$--$1.8\times10^{-12}\eV$ & $(0.60\text{--}3.8)\times10^{-13}\eV$ \\
  $10^{10}\yr$ & $1.4\times10^{-13}$--$2.7\times10^{-12}\eV$ & $(0.35\text{--}5.6)\times10^{-13}\eV$ \\
  \bottomrule
  \end{tabular}
  \end{center}
\item \textbf{2 世代目の較正}: $\mu\approx7\times10^{-13}\eV$ は、$10$--$20\Msun$ の 2 世代目を、遅れが $10^6\yr$ でもすべて回転を失わせる。自由場で \arxiv{2609.02678} の質量を救う道はない。これは同論文が GW241011 について自分で示したこと(E1)と一致する。
\end{itemize}
\end{result}

\begin{caveat}
\begin{itemize}
\item 成長率は解析式で、$\alpha\lesssim0.2$ でしか正確でない。図の上端(大きい $\mu$)の境界は、相対論的な率に替えると動く。
\item 「帯の中のすべての質量」という基準は説明のためのもので、集団解析ではない。
\item 雲の重力波による減衰、伴星による準位の混合、降着は入れていない。どれもスピンを戻さないか(D2)、戻す量が小さい。
\item 2 世代目の自然スピンは 0.69 に固定した。実際は親の質量比とスピンで $0.6$--$0.8$ に広がる。
\end{itemize}
\end{caveat}
